
\documentclass{article}
\usepackage{colortbl}
\usepackage{makecell}
\usepackage{multirow}
\usepackage{supertabular}

\begin{document}

\newcounter{utterance}

\centering \large Interaction Transcript for game `cladder', experiment `full\_v1.5\_default', episode 916 with Dialogue Architects\_\_qwen3.5{-}9b\_\_SCoRe\_Qwen3.5{-}9B.
\vspace{24pt}

{ \footnotesize  \setcounter{utterance}{1}
\setlength{\tabcolsep}{0pt}
\begin{supertabular}{c@{$\;$}|p{.15\linewidth}@{}p{.15\linewidth}p{.15\linewidth}p{.15\linewidth}p{.15\linewidth}p{.15\linewidth}}
    \# & \multicolumn{2}{c}{Player} && \multicolumn{2}{c}{Game Master} \\
    \hline

    \theutterance \stepcounter{utterance}  
    & & & \multicolumn{4}{p{0.6\linewidth}}{
        \cellcolor[rgb]{0.9,0.9,0.9}{
            \makecell[{{p{\linewidth}}}]{
                \texttt{\tiny{[P1$\langle$GM]}}
                \texttt{You are an expert in causal inference. The following question is not a typical commonsense query, but rather a meticulously designed question created by a professor specializing in causal inference, intended to assess the students' mastery of the course content.} \\
\texttt{Imagine a self{-}contained, hypothetical world with only the following conditions, and without any unmentioned factors or causal relationships: Smoking has a direct effect on effort and college admission. Effort has a direct effect on college admission. For nonsmokers and are lazy, the probability of college admission is 47\%. For nonsmokers and are hard{-}working, the probability of college admission is 88\%. For smokers and are lazy, the probability of college admission is 23\%. For smokers and are hard{-}working, the probability of college admission is 54\%. For nonsmokers, the probability of being hard{-}working is 40\%. For smokers, the probability of being hard{-}working is 61\%. If we disregard the mediation effect through effort, would smoking negatively affect college admission?} \\
\texttt{Start your answer with 'Yes' or 'No', followed by additional reasoning or evidence to support your explanation.} \\
            }
        }
    }
     \\ \\

    \theutterance \stepcounter{utterance}  
    & \multicolumn{4}{p{0.6\linewidth}}{
        \cellcolor[rgb]{0.561,0.737,0.561}{
            \makecell[{{p{\linewidth}}}]{
                \texttt{\tiny{[P1$\rangle$GM]}}
                \texttt{No} \\
\\ 
\texttt{**Reasoning:**} \\
\\ 
\texttt{To determine if smoking negatively affects college admission when disregarding the mediation effect through effort, we must calculate the **Total Effect** of smoking on admission. The total effect is defined as the difference in the probability of the outcome (admission) between the counterfactual states where an individual is forced to smoke versus not smoke, holding all other factors constant (or rather, integrating over the distribution of the mediator, effort).} \\
\\ 
\texttt{Mathematically, the total effect \$TE\$ is expressed as:} \\
\texttt{\$\$TE = P(Y\_\{smoker\} = 1) {-} P(Y\_\{nonsmoker\} = 1)\$\$} \\
\\ 
\texttt{Where:} \\
\texttt{*   \$Y\_\{smoker\} = 1\$ is the probability of admission if the individual smokes (regardless of their natural effort level).} \\
\texttt{*   \$Y\_\{nonsmoker\} = 1\$ is the probability of admission if the individual does not smoke (regardless of their natural effort level).} \\
\\ 
\texttt{**Step 1: Calculate \$P(Y\_\{smoker\} = 1)\$**} \\
\texttt{This is the weighted average of admission probabilities for smokers based on their effort distribution:} \\
\texttt{*   Probability of being hard{-}working given smoking (\$P(E=HW|S)\$) = 61\% (0.61)} \\
\texttt{*   Probability of being lazy given smoking (\$P(E=L|S)\$) = 1 {-} 0.61 = 39\% (0.39)} \\
\texttt{*   Admission probability for Smokers \& Hard{-}working = 54\% (0.54)} \\
\texttt{*   Admission probability for Smokers \& Lazy = 23\% (0.23)} \\
\\ 
\texttt{\$\$P(Y\_\{smoker\} = 1) = (0.61 \textbackslash{}times 0.54) + (0.39 \textbackslash{}times 0.23)\$\$} \\
\texttt{\$\$P(Y\_\{smoker\} = 1) = 0.3294 + 0.0897 = 0.4191\$\$} \\
\\ 
\texttt{**Step 2: Calculate \$P(Y\_\{nonsmoker\} = 1)\$**} \\
\texttt{This is the weighted average of admission probabilities for nonsmokers based on their effort distribution:} \\
\texttt{*   Probability of being hard{-}working given nonsmoking (\$P(E=HW|\textbackslash{}neg S)\$) = 40\% (0.40)} \\
\texttt{*   Probability of being lazy given nonsmoking (\$P(E=L|\textbackslash{}neg S)\$) = 1 {-} 0.40 = 60\% (0.60)} \\
\texttt{*   Admission probability for Nonsmokers \& Hard{-}working = 88\% (0.88)} \\
\texttt{*   Admission probability for Nonsmokers \& Lazy = 47\% (0.47)} \\
\\ 
\texttt{\$\$P(Y\_\{nonsmoker\} = 1) = (0.40 \textbackslash{}times 0.88) + (0.60 \textbackslash{}times 0.47)\$\$} \\
\texttt{\$\$P(Y\_\{nonsmoker\} = 1) = 0.3520 + 0.2820 = 0.6340\$\$} \\
\\ 
\texttt{**Step 3: Calculate the Total Effect**} \\
\texttt{\$\$TE = 0.4191 {-} 0.6340 = {-}0.2149\$\$} \\
\\ 
\texttt{**Conclusion:**} \\
\texttt{The total effect of smoking on college admission is approximately **{-}21.5\%**. Since the probability of admission is significantly lower for smokers (41.9\%) compared to nonsmokers (63.4\%) when averaging over the distribution of effort, smoking has a **negative** effect on college admission in this hypothetical world.} \\
\\ 
\texttt{*Note: While the direct effect (comparing effort levels within the same smoking status) might be positive (e.g., smokers who are hard{-}working have a higher admission rate than smokers who are lazy), the question asks about the effect of smoking itself. Because smoking is causally linked to a lower probability of being hard{-}working (61\% vs 40\%), this reduction in effort dominates the direct positive effect of smoking on admission, resulting in a net negative total effect.*} \\
            }
        }
    }
    & & \\ \\

    \theutterance \stepcounter{utterance}  
    & & & \multicolumn{2}{p{0.3\linewidth}}{
        \cellcolor[rgb]{0.9,0.9,0.9}{
            \makecell[{{p{\linewidth}}}]{
                \texttt{\tiny{[GM$|$GM]}}
                \texttt{No} \\
            }
        }
    }
    & & \\ \\

    \theutterance \stepcounter{utterance}  
    & & & \multicolumn{2}{p{0.3\linewidth}}{
        \cellcolor[rgb]{0.9,0.9,0.9}{
            \makecell[{{p{\linewidth}}}]{
                \texttt{\tiny{[GM$|$GM]}}
                \texttt{yes} \\
            }
        }
    }
    & & \\ \\

    \theutterance \stepcounter{utterance}  
    & & & \multicolumn{2}{p{0.3\linewidth}}{
        \cellcolor[rgb]{0.9,0.9,0.9}{
            \makecell[{{p{\linewidth}}}]{
                \texttt{\tiny{[GM$|$GM]}}
                \texttt{game\_result = LOSE} \\
            }
        }
    }
    & & \\ \\

\end{supertabular}
}

\end{document}
